\subsection{Используемые средства разработки}

Разработка интеллектуальной системы анализа обитаемости экзопланет требует тщательного подхода к выбору языков программирования, библиотек визуализации, инструментов анализа данных и протоколов взаимодействия. Учитывая, что основная задача системы заключается в оценке потенциальной обитаемости экзопланет на основе их физических параметров, а также в интерактивной визуализации результатов анализа, необходимо было выбрать такие средства разработки, которые обеспечивали бы:
\begin{itemize}
    \item точность и надежность математических вычислений;
    \item высокую производительность при обработке данных;
    \item качественную визуализацию трехмерных моделей планет;
    \item интерактивное построение графиков и диаграмм;
    \item удобный и отзывчивый пользовательский интерфейс.
\end{itemize}

На этапе предварительного проектирования было проведено сравнительное исследование современных технологий веб-разработки и научных вычислений. В качестве кандидатов были рассмотрены: Python с Django/Flask, JavaScript/Node.js, R с Shiny, MATLAB Web App Server и Julia. Ниже приведён подробный анализ этих технологий с позиции их применимости к разработке системы анализа экзопланет.

Python — высокоуровневый язык программирования с богатой экосистемой научных библиотек. Он активно используется в астрономии, анализе данных и машинном обучении. Python предоставляет множество специализированных библиотек для научных вычислений, таких как NumPy, SciPy и Pandas, которые идеально подходят для обработки астрономических данных и выполнения сложных математических расчетов.

Для веб-разработки на Python наиболее часто используются фреймворки Flask и Django. В данном проекте был выбран Flask (версия 2.3.3) как более легковесное и гибкое решение, позволяющее создавать REST API для взаимодействия с клиентской частью приложения. Flask обеспечивает простую интеграцию с научными библиотеками Python и предоставляет все необходимые инструменты для разработки веб-сервисов.

JavaScript, будучи основным языком веб-браузеров, был выбран для реализации клиентской части приложения. Современный JavaScript (ES6+) предоставляет мощные возможности для создания интерактивных пользовательских интерфейсов и визуализации данных. Особенно важными для проекта стали следующие библиотеки и фреймворки:

\begin{itemize}
    \item Three.js (версия 0.132.2) — мощная библиотека для создания трехмерной графики в веб-браузере, которая используется для визуализации моделей экзопланет с учетом их физических характеристик;
    \item Plotly.js (версия 2.27.1) — библиотека для создания интерактивных научных графиков и диаграмм, позволяющая наглядно представлять результаты анализа;
    \item Tippy.js (версия 6) — инструмент для создания информативных всплывающих подсказок, улучшающих пользовательский опыт при работе с системой.
\end{itemize}

Для обработки данных и научных вычислений были выбраны следующие Python-библиотеки:

\begin{itemize}
    \item pandas (версия 2.1.0) — библиотека для эффективной работы с табличными данными, используемая для обработки каталогов экзопланет;
    \item NumPy (версия 1.25.2) — фундаментальный пакет для научных вычислений, применяемый для расчета физических параметров и индексов обитаемости;
    \item scikit-learn (версия 1.3.0) — библиотека машинного обучения, используемая для анализа и классификации экзопланет по их потенциальной обитаемости.
\end{itemize}

Выбранный стек технологий обеспечивает оптимальный баланс между производительностью научных вычислений на серверной стороне и интерактивностью пользовательского интерфейса на клиентской стороне. Python с его научными библиотеками предоставляет надежную основу для математических расчетов и анализа данных, в то время как JavaScript с современными фреймворками обеспечивает богатый пользовательский интерфейс и качественную визуализацию результатов.

JavaScript в контексте веб-браузера отлично подходит для обработки интерактивных элементов системы анализа экзопланет, поскольку работает по асинхронной, событийно-ориентированной модели, что особенно важно при визуализации трехмерных моделей планет, построении графиков и динамической подгрузке данных анализа.

JavaScript активно используется на стороне клиента для отображения интерфейса системы, в том числе для:
\begin{itemize}
    \item обработки пользовательского ввода параметров экзопланет;
    \item асинхронной загрузки и отображения результатов анализа;
    \item управления трехмерными моделями планет через Three.js;
    \item визуализации научных данных через Plotly.js.
\end{itemize}

Кроме того, использование современных веб-технологий позволяет эффективно работать с WebGL для 3D-рендеринга, Canvas API для построения графиков, и асинхронными вычислениями для анализа данных — все эти задачи легко реализуются средствами JavaScript.

В рамках проекта JavaScript был выбран как основной язык реализации клиентской части системы в силу своей гибкости, широкой поддержки визуализации данных и отличной интеграции с современными веб-технологиями.

R с Shiny — это мощная комбинация для создания интерактивных веб-приложений с научными вычислениями. R традиционно используется в статистике и анализе данных, а Shiny предоставляет фреймворк для создания веб-интерфейсов к R-скриптам.

Shiny позволяет быстро создавать интерактивные дашборды и визуализации данных, что могло бы быть полезно для отображения статистики экзопланет. Однако у этого стека есть ряд ограничений:
\begin{itemize}
    \item сложность интеграции с современными веб-технологиями;
    \item ограниченные возможности кастомизации интерфейса;
    \item меньшая производительность при работе с 3D-графикой.
\end{itemize}

В связи с этим, R с Shiny не был выбран для данного проекта, несмотря на его сильные стороны в области статистического анализа.

MATLAB Web App Server предоставляет возможность создавать веб-приложения на основе MATLAB-скриптов. MATLAB имеет мощные средства для научных вычислений и обработки данных, что потенциально полезно для анализа параметров экзопланет.

Однако его использование связано с рядом ограничений:
\begin{itemize}
    \item высокая стоимость лицензии;
    \item сложность интеграции с современными веб-технологиями;
    \item ограниченные возможности создания интерактивных интерфейсов.
\end{itemize}

Julia — относительно новый язык программирования, специально разработанный для научных вычислений. Он обеспечивает производительность, сравнимую с C, при синтаксисе, близком к Python. Julia имеет растущую экосистему пакетов для астрономических вычислений.

Однако экосистема Julia для веб-разработки пока недостаточно развита, что затрудняет создание современных интерактивных веб-приложений. Кроме того, сообщество Julia меньше, чем у Python или JavaScript, что может затруднить поиск решений и готовых компонентов.

Проведённый анализ современных языков программирования и сопутствующих технологий позволил обоснованно выбрать оптимальные средства разработки для реализации системы анализа обитаемости экзопланет. С учётом требований к точности вычислений, визуализации данных, интерактивности интерфейса и масштабируемости, было принято решение использовать комбинацию Python и JavaScript.

Python с его научными библиотеками обеспечивает надежную основу для математических расчетов и анализа данных, в то время как JavaScript с современными фреймворками предоставляет богатые возможности для создания интерактивного пользовательского интерфейса и визуализации результатов. Такой подход позволяет эффективно решать как задачи научных вычислений, так и задачи создания современного веб-интерфейса. 